% SI Text 5: extended results. Full treatment of results the main text states in one
% sentence; drawn from the long version's Sections 4 to 7 (paragraphs and footnotes),
% with Figs. S1 to S4.
\section{Extended results}\label{si:extended}

This section gives the full treatment of results that the main text states in one sentence. Each subsection opens with the main-text statement it supports.

\subsection*{The value of targeting decreases in the budget}
\emph{Main text:} ``The value of targeting, the expected output the optimal allocation adds over equal division, decreases in the budget and vanishes as the budget grows.''

How much targeting matters depends on the budget. A researcher's expected output increases with every dollar of resources, but with diminishing marginal returns, and it never exceeds $AK_i$: expected output is bounded above by a quantity set by capability, which it approaches but does not attain. When the budget is large, optimal and uniform funding therefore differ little in the output they produce: under either, every researcher's output is near its upper bound, and the difference between the two vanishes as the budget increases (Corollary~\ref{cor:targeting-vanishes}, \ref{app:gap-rule}). When the budget is small, dollars are scarce, researchers differ in the marginal value of a dollar, and output depends on where each dollar goes. Targeting matters where the budget is small relative to researchers' resources and little where the budget is large (Fig.~\ref{fig:targeting-value}). This dependence on the budget recurs later: it sets how much output is lost when seed grants and lotteries override the optimal allocation, and it determines when optimal funding concentrates rather than spreads (both in the main text).

% Figure 2: the value of targeting decreases in the budget.
% Computed exactly from Proposition 1 (no simulation): population of n = 75 with
% K, R0 ~ Pareto(alpha = 2, min 1), rho = 0, A = 1, seed 659 (numpy default_rng);
% value of targeting = (f_opt - f_unif) / (f_unif - f_0), in percent;
% b = budget / total baseline resources. Data: fig2-curve.dat, generated by
% verify_gap_rule.py's companion snippet (see state.md, 2026-09-05).
\begin{figure}[tp]
\centering
\begin{tikzpicture}
\begin{axis}[
  width=0.82\textwidth, height=5.4cm,
  axis lines*=left,
  xmin=0, xmax=2, ymin=0, ymax=175,
  xtick={0, 0.5, 1, 1.5, 2},
  ytick={0, 50, 100, 150},
  xlabel={budget scale $b$},
  ylabel={value of targeting (\%)},
  ylabel style={font=\small}, xlabel style={font=\small},
  tick label style={font=\small},
  axis line style={line width=0.4pt},
  tick style={line width=0.4pt, black},
  clip=false,
]
\addplot[black, line width=0.9pt, smooth] table {fig2-curve.dat};
\end{axis}
\end{tikzpicture}
\caption{The value of targeting against the budget. The optimal allocation's gain over uniform funding, as a fraction of uniform funding's gain over no funding, as the budget scale $b$ grows, for one population drawn from the default distributions. Corollary~\ref{cor:targeting-vanishes} proves that the value vanishes as the budget grows; that it falls throughout is a numerical finding (\ref{app:gap-rule}).}
\label{fig:targeting-value}
\end{figure}


\subsection*{The value of targeting relative to uniform funding}
\emph{Main text:} ``For a funder allocating with records and review, targeting adds more research output than uniform funding's entire gain over no funding in a heavy-tailed field with reliable review and a tight budget, and about a quarter of that gain in the default field at the largest budget we sweep.''

The review-informed funder's gain over uniform funding, as a fraction of uniform funding's gain over no funding, decreases from 1.17 to 0.72 across budget scales $b = 0.2$ to $2$ with heavy-tailed capability ($\alpha_K = 1.3$) and reliable review ($\tau_K = 0.3$), and from 0.74 to 0.25 in the default field with the default signal ($\alpha_K = 2$, $\tau_K = 1$). Two rounds; 200 simulated populations per point (the populations of Fig.~4\textit{A}). Giving up targeting loses little output under the same conditions that make review worth little: where capability is evenly spread, the budget is large, or review is uninformative.

\subsection*{Overtrusting review}
\emph{Main text:} ``A funder that overtrusts review, treating a noisy signal as reliable, can do worse than one that ignores review altogether, and overtrusting review loses more output than undertrusting it.''

One condition qualifies these results: the funder must weight review at its true informativeness. A funder that overtrusts review, treating a noisy signal as reliable, can do worse than one that ignores review altogether: with moderately spread capability, overtrust turns review's value negative; with heavy-tailed capability it reduces the value without reversing it (Fig.~\ref{fig:overtrust}).\footnote{A funder that overtrusts tenfold (true $\tau_K = 3$, belief $0.3$) loses more than review's full value at the default capability distribution; a funder that undertrusts tenfold (true $\tau_K = 0.3$, belief $3$) retains over half of it.} Overtrusting review loses more output than undertrusting it.

% Figure 6 (S6): the value of review under misweighted trust.
% y = review's value (S5 - S4) as percent of no-funding output; x = the funder's
% believed noise (log scale); true noise tau_K = 3 (dotted rule). Two capability
% distributions. Error bars: one standard error, 200 seeds per cell.
% Data: fig6-overtrust.dat, from sweep_results/D_misspecified_trust (hash 21b0d9a).
\begin{figure}[tp]
\centering
\begin{tikzpicture}
\begin{axis}[
  width=0.82\textwidth, height=6.0cm,
  axis lines*=left,
  xmode=log,
  xmin=0.1, xmax=10, ymin=-5, ymax=45,
  xtick={0.1, 0.3, 1, 3, 10},
  xticklabels={0.1, 0.3, 1, 3, 10},
  ytick={0, 10, 20, 30, 40},
  xlabel={believed noise $\tau_K$},
  ylabel={value of review (\% of the gain from funding)},
  ylabel style={font=\small}, xlabel style={font=\small},
  tick label style={font=\small},
  axis line style={line width=0.4pt},
  tick style={line width=0.4pt, black},
  clip=false,
]
\addplot[gray, dashed, line width=0.4pt, domain=0.1:10] {0};
\draw[gray, dotted, line width=0.5pt] (axis cs:3,-2) -- (axis cs:3,20);
\node[font=\small, gray, anchor=west] at (axis cs:3.15,18) {true noise};
\addplot[black, line width=0.7pt, error bars/.cd, y dir=both, y explicit]
  table[x=belief, y=p13, y error=s13] {fig6-overtrust.dat};
\addplot[black, line width=0.7pt, error bars/.cd, y dir=both, y explicit]
  table[x=belief, y=p20, y error=s20] {fig6-overtrust.dat};
\node[font=\small, anchor=west] at (axis cs:11,25.5) {$\alpha_K = 1.3$};
\node[font=\small, anchor=west] at (axis cs:11,3.9)  {$\alpha_K = 2$};
\end{axis}
\end{tikzpicture}
\caption{Overtrusting review loses more research output than undertrusting it. The value of review as a percent of the research output that a correctly calibrated review-informed funder adds over no funding, as the funder's believed signal noise varies while the true noise is $\tau_K = 3$ (dotted rule); bars are one standard error. At the default capability distribution, overtrust takes the value below zero. Settings: \ref{app:simulation}.}
\label{fig:overtrust}
\end{figure}


\subsection*{Grant size, learning, and the schedule}
\emph{Main text:} ``The shift of spending toward later rounds comes from the output researchers produce whether or not they are funded, which compounds and reveals capability while the funder waits.''

How much a funder can learn from its own grants depends on their size. A researcher on a small grant produces output nearly proportional to the grant and nearly independent of capability, so a funder whose grants are small learns almost nothing it can exploit later: in a field whose researchers have few resources, with no review signal, such a funder gains almost nothing over uniform funding. Large grants change this. As grants approach the scale of capability itself, output separates the capable from the rest, and the funder recovers most of the value of discrimination, on a schedule that spends a small share early, so that output can be observed, and the rest once the funder has learned who is capable (Fig.~\ref{fig:depth-schedule}).\footnote{In a community whose researchers have few resources, with heavy-tailed capability and no review signal, over six rounds, the Bayesian funder's gain over uniform funding increases from a third of one percent of uniform funding's output at the default grant size to nine percent at six times that size and ten percent at twelve times, roughly the gain a reliable review signal delivers at the default size.}

The shift of spending toward later rounds comes from the output researchers produce whether or not they are funded: that output compounds and reveals capability while the funder waits. Growth from grant-added output alone would shift the schedule slightly toward earlier rounds, since earlier grants have longer to compound.\footnote{With growth from unfunded output switched off, the schedule's center of mass is 0.48 to 0.49, below even; with growth from grant-added output switched off, 0.53 to 0.68 (five rounds).}

% Figure (S7): the optimal schedule at standard and deep funding.
% y = share of the budget spent in each round by the forward funder (S8); T = 6;
% poverty corner (r_min = 0.001), heavy tail (alpha_K = 1.3), no review signal
% (tau_K = 100), negligible compounding (eps = 1e-4), budget scaled against
% capability (budget_ref = "K"). b = 0.5: 50 seeds; b = 3: 24 seeds (max SE 0.006);
% round-1 share at b = 3 matches the canonical 200-seed run (0.104 vs 0.110).
% Data: fig7-schedule.dat, generated by for-claude/verify_s7.R (schedule addendum).
\begin{figure}[tp]
\centering
\begin{tikzpicture}
\begin{axis}[
  width=0.66\textwidth, height=5.4cm,
  axis lines*=left,
  xmin=1, xmax=6, ymin=0.08, ymax=0.20,
  xtick={1, 2, 3, 4, 5, 6},
  ytick={0.10, 0.15, 0.167, 0.20},
  yticklabels={0.10, 0.15, even, 0.20},
  xlabel={round},
  ylabel={share of budget},
  ylabel style={font=\small}, xlabel style={font=\small},
  tick label style={font=\small},
  axis line style={line width=0.4pt},
  tick style={line width=0.4pt, black},
  clip=false,
]
\addplot[black, line width=0.7pt, mark=*, mark size=1.2pt] table[x=round, y=b05] {fig7-schedule.dat};
\addplot[black, line width=0.7pt, mark=o, mark size=1.4pt] table[x=round, y=b3] {fig7-schedule.dat};
\node[font=\small, anchor=west] at (axis cs:6.15,0.167) {default grant size ($b = 1$)};
\node[font=\small, anchor=west] at (axis cs:6.15,0.183) {large grants ($b = 6$)};
\end{axis}
\end{tikzpicture}
\caption{The forward-looking funder's spending by round at two grant sizes. The share of the budget spent in each of six rounds, in a field with heavy-tailed capability, few baseline resources, and no review, at the default grant size (filled) and at six times that size (open). Settings: \ref{app:simulation}.}
\label{fig:depth-schedule}
\end{figure}


\subsection*{Concentration and the production function}
\emph{Main text:} ``Stronger complementarity between capability and resources concentrates the review-informed funder's grants when the budget is tight in a heavy-tailed field, spreads them when the budget is ample, and has a non-monotone effect where capability is evenly spread.''

Seed grants and lotteries both spread funds more widely than targeting does, and both raise the older question of whether funding should be concentrated among a few researchers or spread across many. One might hope that the production function settles this question: the harder it is to substitute resources for capability, the more the money should go to the few researchers who can use it. To test this, we embed our production function in a family indexed by an exponent $\gamma$, from Cobb-Douglas ($\gamma = 0$), under which a shortfall in either input can be offset by more of the other, through our form ($\gamma = -1$), to Leontief, under which output depends only on the scarcer input; the harder it is to substitute one input for the other, the stronger the complementarity (\ref{app:technology}). The production function does not settle the question (Fig.~\ref{fig:concentration}). With a tight budget in a heavy-tailed field, stronger complementarity does concentrate the review-informed funder's grants. With an ample budget, stronger complementarity spreads them, a pattern consistent with a funder that, once the budget allows, brings every researcher's resources up toward the level their capability can use. Where capability is evenly spread and the budget ample, the relation is not even monotone: concentration peaks at an intermediate complementarity and declines toward both ends.\footnote{Gini coefficients of the review-informed funder's grants, from Cobb-Douglas to Leontief: 0.92 to 0.95 in a heavy-tailed field with a tight budget; 0.35 to 0.21 in the default field with an ample budget; 0.63 to 0.45 in a heavy-tailed field with an ample budget. In the evenly spread field with an ample budget the Gini coefficient is 0.16 at Cobb-Douglas, peaks at 0.196 near $\gamma = -6$, and falls to 0.185 at $\gamma = -12$ (3.5 standard errors above the $\gamma = -12$ level).} How widely funds should be spread is set jointly by the production function, the budget, and the field's capability distribution; no one of the three decides it alone.

% Figure (S8): how concentrated the informed allocation is, across the family.
% y = Gini coefficient of the review-informed funder's grants (single round);
% x = CES exponent gamma (0 = Cobb-Douglas, -1 = our harmonic form, limit =
% Leontief, plotted as detached marks at the left edge). Three field-budget pairs.
% Ground: sigma_tierB_concentration (+ leontief endpoint, + 400-seed refinement for
% the evenly spread curve), 200 seeds per cell, re-derived in-session (verify_s9.R).
% Leontief cells carry the kink caveat (ties broken by fill order, n_steps = 800).
\begin{figure}[tp]
\centering
\begin{tikzpicture}
\begin{axis}[
  width=0.82\textwidth, height=6.0cm,
  axis lines*=left,
  xmin=-14.8, xmax=0, ymin=-0.16, ymax=0.06,
  xtick={-14, -12, -6, -3, -1, 0},
  xticklabels={Leontief, $-12$, $-6$, $-3$, $-1$, $0$},
  ytick={-0.15, -0.10, -0.05, 0, 0.05},
  yticklabels={$-0.15$, $-0.10$, $-0.05$, 0, 0.05},
  xlabel={CES exponent $\gamma$ (complements $\leftarrow$, substitutable $\rightarrow$)},
  ylabel={change in grant concentration (Gini)},
  ylabel style={font=\small}, xlabel style={font=\small},
  tick label style={font=\small},
  axis line style={line width=0.4pt},
  tick style={line width=0.4pt, black},
  clip=false,
]
\draw[gray, dotted, line width=0.5pt] (axis cs:-1,-0.16) -- (axis cs:-1,0.06);
\node[font=\scriptsize, gray, anchor=south] at (axis cs:-1,0.061) {our form};
\addplot[black, line width=0.7pt, mark=*, mark size=1.1pt] table[x=g, y=dgini] {fig13-heavy01.dat};
\addplot[black, line width=0.7pt, mark=o, mark size=1.3pt] table[x=g, y=dgini] {fig13-def1.dat};
\addplot[black, line width=0.7pt, mark=triangle*, mark size=1.4pt] table[x=g, y=dgini] {fig13-even1.dat};
\addplot[black, only marks, mark=*, mark size=1.1pt] coordinates {(-14,0.031)};
\addplot[black, only marks, mark=o, mark size=1.3pt] coordinates {(-14,-0.147)};
\addplot[black, only marks, mark=triangle*, mark size=1.4pt] coordinates {(-14,0.009)};
\node[font=\scriptsize, anchor=north] at (axis cs:-8.5,0.014) {heavy-tailed, tight};
\node[font=\scriptsize, anchor=south] at (axis cs:-10.5,-0.047) {default, ample};
\node[font=\scriptsize, anchor=south] at (axis cs:-7.5,0.042) {evenly spread, ample};
\end{axis}
\end{tikzpicture}
\caption{Complementarity concentrates funding only when the budget is small. The change in the Gini coefficient of the review-informed funder's grants, relative to its Cobb-Douglas level, as the production function moves from Cobb-Douglas ($\gamma = 0$) toward Leontief (detached marks at the left edge), for three field-and-budget combinations. Settings: \ref{app:simulation}.}
\label{fig:concentration}
\end{figure}

